\section{The anatomy layer: the honest state}\label{sec:typeII}

\textbf{[REWRITTEN after adversarial verification, June 10 (late).]}
An earlier draft of this section claimed the anatomy layer glues from
classical inputs. Verification showed the claim rested on a
band-arithmetic error, and the correct picture is structural:
\emph{in the band, $u+u'>\tfrac23$ and $u+2u'>1$ always}. Hence
$R=x^{1-u-u'}<x^{1/3}=y$: every prime $\ell>y$ exceeds the progression
length, the ``small-$\Lambda$'' regime is empty, and the subtraction
term of Proposition~\ref{prop:species} is exactly
\[
\Sigma(\theta)\;=\;\sum_{t\le R}\mathbf 1\big[P^{+}(a_0+qt)>y\big]\,
\e{t\theta}
\]
--- a genuine arithmetic exponential sum with no per-term damping.
Moreover every weighted sub-object here lives in progressions whose
modulus exceeds the square root of the sum length ($q>\sqrt{A/b}$
always, since $u+2u'>1$): beyond Vaughan, beyond
Bombieri--Vinogradov, at the dispersion frontier. Cauchy--Schwarz
routes are structurally incapable of recovering the $1/q$ congruence
density (even perfect square-root cancellation gives
$\sqrt A>R$). The unweighted digit layer
(\S\S\ref{sec:reduction}--\ref{sec:major}) is immune because its
progression sums are exact geometric series; the anatomy weight is
where that exactness ends.

\begin{hypothesis}[E3; anatomy dispersion]\label{hyp:E3}
\statusO{} (structural; the residual core of D$^\dagger$-anatomy)\quad
Uniformly in $0<|\lambda|\le\Llog^{C}$,
\[
\frac{1}{\pi(Q)}\sum_{q\sim Q}\ \sum_{\substack{p\sim P\\
p\ \mathrm{minor}}}\ \sum_{0<\mu<p}\frac{1}{\langle\mu\rangle}
\Big|\sum_{t\le R}\mathbf 1\big[P^{+}(a_0+qt)>y\big]\,\e{t\theta_\mu}
\Big|\;\ll\;\pi(P)\,R\,\Llog^{-c-1}.
\]
\end{hypothesis}

\begin{remark}[evidence and attack surface]
The empirical record bears directly on E3: the WP1.5/WP2.4 probes
measured the \emph{fully decorated} sums (anatomy weight included) and
found square-root cancellation at every stratum
(Appendix~\ref{app:numerics}); the fine-class probes confirm it within
classes. The plausible tools are dispersion in $q$ (the phase
differences $\theta_{\mu_1}-\theta_{\mu_2}=(\mu_1-\mu_2)q/p$ move with
$q$; BFI-style), the deep large sieve as the unconditional backstop,
and the friable-in-APs literature (Fouvry--Tenenbaum) for the weight's
own equidistribution. A pure $L^2$ in $t$ fails (the weight has
variance $\asymp R$); the $\mu$- and $q$-averages are where the
cancellation must come from. For the \emph{lower-bound track}, sieve
(Buchstab) decompositions of the friability indicator with positivity
may substitute for E3 on a positive-proportion subclass --- to be
investigated.
\end{remark}

\begin{remark}[what survives from the earlier draft]
The bilinear computation (Cauchy--Schwarz with the congruence opened
by additive characters) is algebraically valid and recorded in the
repository history; it is not budget-compliant for the reason above.
The linear $\ell$-resonance count survives as a tool with corrected
hypotheses (grid-site counting; it needs gap control \emph{at the
scale used}, the same all-scale subtlety as
Remark~\ref{rem:M2-strengthening}). The extreme range
($b\le x^{\eta_1}$) is primes in APs with modulus beyond the square
root of the length --- the BFI/large-moduli frontier, not
``standard toolbox.''
\end{remark}
